% Full proof and printed finite witnesses for the companion supplement.
% Requires amsmath, amssymb, amsthm, booktabs, xurl, hyperref.
\section{Exact stability at seven and all cyclotomic characters}
\label{seven:section}

Fukuda's stability theorem reduces a class-group calculation in an entire
$\mathbb Z_p$-extension to two successive finite layers
\cite[Theorem~1(1), p.~264]{sevenFukuda1994}.
We give its short finite-level proof in the one-prime situation needed here
and certify the two arithmetic inputs. This gives an exact finite Iwasawa
module, including its finite submodule, and hence determines every strict
cyclotomic specialization.

Put
\[
 F=\mathbb Q(\alpha),\quad \alpha^3=269,\qquad
 K=\mathbb Q(\omega),\quad E=KF,
\]
where $\alpha$ is the real root and $\omega$ is a primitive cube root of unity.
Let $Q_n$ be the degree-$7^n$ layer of the cyclotomic
$\mathbb Z_7$-extension of $\mathbb Q$, and set
$F_n=FQ_n$, $K_n=KQ_n$, and $E_n=EQ_n$.
Write $A(k)=\operatorname{Cl}(k)\{7\}$, and
\[
 \Delta=\operatorname{Gal}(E/K)=\langle\delta\rangle,\qquad
 \Gamma=\operatorname{Gal}(E_\infty/E).
\]
Both faithful characters of $\Delta$ take values in $\mu_3\subset\mathbb Z_7$.
Choose $\gamma\in\Gamma$ with
$\langle\varepsilon_7\rangle(\gamma)=8$, and put
$T=\gamma-1$ and $\Lambda=\mathbb Z_7[[T]]$.

For a finite extension $V/\mathbb Q_7$, with integer ring $\mathcal O$, and
a continuous character $\kappa:\Gamma\longrightarrow\mathcal O^\times$,
let
\[
 D_{\chi,\kappa}=(V/\mathcal O)(\chi\kappa).
\]
The characters here are inflated through
$\operatorname{Gal}(E_\infty/K)=\Delta\times\Gamma$.
The strict Selmer group means
\[
 \operatorname{Sel}_{\rm str}(K,D)=
 \ker\left(
 H^1(K,D)\longrightarrow
 \prod_{v\mid7}H^1(K_v,D)
 \times\prod_{v\nmid7}H^1(I_v,D)\right).
 \tag{S1}\label{seven:strict-definition}
\]
The second product is over finite places; imposing the usual condition at
infinite places has no effect on seven-primary cohomology.

\begin{theorem}\label{seven:main}
For every $n\ge0$,
\[
 A(F_n)\simeq\mathbb Z/7\mathbb Z.
 \tag{S2}\label{seven:class-stability}
\]
All norms between these groups are isomorphisms and all cyclotomic actions
are trivial. For either faithful $\chi$ of $\Delta$,
\[
 \left(\varprojlim_n A(E_n)\right)^\chi
 \simeq\Lambda/(7,T).
 \tag{S3}\label{seven:iwasawa}
\]
For every continuous $\kappa$ as above, with $u=\kappa(\gamma)$,
\[
 \operatorname{Sel}_{\rm str}(K,D_{\chi,\kappa})
 \simeq\mathcal O/(7,u-1).
 \tag{S4}\label{seven:all-characters}
\]
These assertions are unconditional.
\end{theorem}

The finite certificate in this proof uses exact algebraic integers, a
matrix of residue characters, and bounded convergent series. A list of
units alone does not certify completeness of a class group or a unit
basis. The certificate proves precisely the seven-part required in
\eqref{seven:class-stability}; it does not need a determination of every
prime divisor of either class number.

\section*{Proof overview}

We first compute the quotient of each class number by the index of an
explicit group of units. A universal regulator lower bound makes the
possible denominator finite. Rigorous intervals isolate the quotient as
the integer seven, while independent residue characters prove that the
unit index is prime to seven. In the cyclotomic tower of $F$, exactly one
prime ramifies. Class field theory therefore identifies class-group
coinvariants with the lower class group. Equality at the first two levels
and Nakayama's lemma force stability at every level. The two faithful
cubic-character parts over $E$ are recovered from the conjugation-fixed
part. Finally, inflation and restriction preserve both strict local
conditions because the faithful cubic character has no invariants.
The resulting specialization of the actual finite module is
\eqref{seven:all-characters}.

\subsection{A finite certification principle}
\label{seven:finite-principle}

\begin{lemma}\label{seven:isolation}
Let $k$ be a number field with regulator $R_k>1/4$. Let $U_0$ contain all
roots of unity of $k$ and be generated modulo torsion by $r$ independent
units, where $r$ is the unit rank. Put
$m=[\mathcal O_k^\times:U_0]$ and $R_0=mR_k$.
Suppose $B$ is an integer with $m<B$ and an interval calculation proves
\[
 \left|\frac{h_k}{m}-7\right|<\frac1B.
 \tag{S5}\label{seven:isolation-bound}
\]
Then $h_k=7m$. If $r$ homomorphisms
$\mathcal O_k^\times\longrightarrow\mathbb F_7$ have a nonsingular
evaluation matrix on the specified units, then $7\nmid m$, and
$A(k)\simeq\mathbb Z/7\mathbb Z$.
\end{lemma}

\begin{proof}
The relation $R_0=mR_k$ follows by comparing the covolumes of the logarithmic
unit lattices. Since $h_k-7m$ is an integer of absolute value less than
$m/B<1$, it vanishes. If $7$ divided $m$, the lattice inclusion
$U_0/\mu(k)\subset\mathcal O_k^\times/\mu(k)$ would have a singular
reduction modulo seven. Equivalently, there would be a product of the
specified units, with exponents not all divisible by seven, which is a
seventh power of a unit. Every displayed homomorphism would kill this
product, contradicting the nonsingular matrix. Thus
$v_7(h_k)=1$, as asserted.
\end{proof}

The regulator bound used here is Friedman's
\cite[Theorem~B, pp.~599--600]{sevenFriedman1989}: $R_k>1/4$ outside three
exceptional totally complex sextic fields. Both fields used below have real
embeddings and degrees three and twenty-one, so this exception is irrelevant.
In particular, one may take $B=\lceil4R_0\rceil$, with an upper enclosure of
$R_0$ used when calculating the ceiling.

For a rational prime $\ell\equiv1\pmod7$ and a prime
$\mathfrak l$ of $k$ with residue field $\mathbb F_\ell$, choose
$z\in\mathbb F_\ell^\times$ of order seven. The character
\[
 u^{(\ell-1)/7}\equiv z^{\,\ell_{\mathfrak l}(u)}
 \pmod{\mathfrak l},\qquad
 \ell_{\mathfrak l}(u)\in\mathbb F_7
 \tag{S6}\label{seven:residue-character}
\]
is a homomorphism on all units. This supplies the homomorphisms in the lemma.
The fields here have a real embedding and odd degree, so their roots of
unity are exactly $\{\pm1\}$.

\subsection{The zeta residues as bounded finite sums}
\label{seven:analytic}

Set $N=217083=3\cdot269^2$, and let $\rho$ be the standard two-dimensional
representation of $\operatorname{Gal}(E/\mathbb Q)\simeq S_3$. Then
\[
 L_0(s)=L(s,\rho)=\frac{\zeta_F(s)}{\zeta(s)},\qquad
 D_F=-N.
 \tag{S7}\label{seven:base-L}
\]
For completeness, the discriminant exponent at $269$ is two by tame total
ramification. At three, $v_3((-1)^3-269)=3>2v_3(3)=2$, so Hensel's
lemma gives a root in $\mathbb Q_3$ and
$F\otimes\mathbb Q_3=\mathbb Q_3\times\mathbb Q_3(\omega)$.
The exponent at three is therefore one. These observations also identify
the ramified Euler factors below.

Let $\psi$ be the order-seven character modulo $49$ with
$\psi(3)=\zeta_7$, and put $L_j=L(s,\rho\otimes\psi^j)$.
For $1\le j\le6$, the character $\psi^j$ is even and primitive of conductor
$49$: $3$ has order $42$ modulo $49$, and $-1=3^{21}$.
Its conductor cannot divide seven because
$|(\mathbb Z/7\mathbb Z)^\times|=6$.
Since $\rho$ is unramified at seven, twisting its two unramified dimensions
gives conductor exponent four there. Thus
\[
 N_0=N,\quad N_j=49^2N\ (j\ne0),\qquad
 \Lambda_j(s)=N_j^{s/2}\Gamma_{\mathbb C}(s)L_j(s),
 \quad
 \Gamma_{\mathbb C}(s)=2(2\pi)^{-s}\Gamma(s).
 \tag{S8}\label{seven:conductors}
\]
The Dedekind functional equation gives
$\Lambda_0(s)=\Lambda_0(1-s)$
\cite[Remark~18.14]{sevenSutherland2015}.
The representation is induced from a nontrivial cubic Hecke character of
$K$ of conductor $(269)$, whose conjugate is its distinct inverse.
Hecke's theta construction therefore gives the cuspidal weight-one form
$f(z)=\sum_{n\ge1}a_ne^{2\pi inz}$ of level $N$ and nebentypus
$\eta=\left(\frac{-3}{\,\cdot\,}\right)$; see
\cite[Theorem~5]{sevenSimard2015} for the theta transformation and
\cite[\S3]{sevenDarmonLauderRotger2015} for the weight-one induced
representation and cuspidality criterion.

Here is a derivation of the twisting sign, to fix its orientation.
For a primitive even character $\xi\bmod M$, $M=49$, set
$\tau(\xi)=\sum_{a\bmod M}\xi(a)e^{2\pi ia/M}$ and
\[
 A_a=\begin{pmatrix}1&a/M\\0&1\end{pmatrix},\qquad
 W_Q=\begin{pmatrix}0&-1\\Q&0\end{pmatrix}.
\]
With the normalized weight-one slash operator, finite Fourier inversion gives
\[ f_\xi=\tau(\bar\xi)^{-1}\sum_a\bar\xi(a)f|A_a. \]
Choose $b$ with $Nab\equiv-1\pmod M$. Direct multiplication gives
\[
 A_aW_{NM^2}=
 \begin{pmatrix}(1+Nab)/M&a\\Nb&M\end{pmatrix}M W_NA_b.
 \tag{S9}\label{seven:fricke-matrix}
\]
The first matrix has determinant one, lies in $\Gamma_0(N)$, and has
lower-right entry $M$, with $\eta(M)=1$.
Also $\bar\xi(a)=\xi(N)\xi(b)$.
The base functional equation gives $f|W_N=-if$ by Mellin inversion.
Using $\tau(\xi)\tau(\bar\xi)=M$ consequently gives
\[
 \Lambda_j(s)=\varepsilon_j\Lambda_{7-j}(1-s),\qquad
 \varepsilon_j=\psi^j(N)\frac{\tau(\psi^j)^2}{49}.
 \tag{S10}\label{seven:epsilon}
\]
There is a short exact verification of the Gauss sums.
The relation $3^6=1-7\bmod49$ gives
$\psi(1+7b)=\zeta_7^b$. Grouping terms as $a+7t$ shows that only
$a\equiv-j\pmod7$ survives in $\tau(\psi^j)$.
Since $N\equiv3^{33}\pmod{49}$, the results are
\[
\begin{array}{c|rrrrrr}
 j&1&2&3&4&5&6\\ \hline
 \tau(\psi^j)/7=\zeta_{49}^{\,b_j}
   &48&19&18&31&30&1\\
 \varepsilon_j=\zeta_{49}^{\,c_j}
   &33&10&43&6&39&16
\end{array}
 \tag{S11}\label{seven:gauss-table}
\]
where the entries in the last two rows are the exponents $b_j,c_j$.

The coefficients $a_n$ are multiplicative and have the following complete
prime-power specification, including $a_{p^0}=1$:
\[
\begin{array}{c|l}
 p& a_{p^k}\quad(k\ge1)\\ \hline
269&0\\
3&1\\
p\equiv2\pmod3,\ p\ne269
   &1\text{ for even }k,\quad0\text{ for odd }k\\
p\equiv1\pmod3,\ 269^{(p-1)/3}\equiv1\pmod p&k+1\\
p\equiv1\pmod3,\ 269^{(p-1)/3}\not\equiv1\pmod p
   &1,-1,0\text{ for }k\equiv0,1,2\pmod3 .
\end{array}
 \tag{S12}\label{seven:coefficients}
\]
These are the Euler factors for the identity, transposition, and
three-cycle classes, together with the inertia invariants at the two
ramified primes. In particular, $|a_n|\le d(n)\le n$.

Put $\lambda_j=2\pi/\sqrt{N_j}$.
Splitting the Mellin integral at one and using
\eqref{seven:epsilon} yields, with
$E_1(x)=\int_x^\infty e^{-t}\,dt/t$,
\[
 L_j(1)=
 \sum_{n\ge1}\frac{a_n\psi^j(n)e^{-\lambda_jn}}{n}
 +\varepsilon_j \lambda_j
 \sum_{n\ge1}a_n\overline{\psi^j(n)}E_1(\lambda_jn).
 \tag{S13}\label{seven:L1-formula}
\]
For $j=0$ the character in this formula is identically one, including at
multiples of seven, and $\varepsilon_0=1$.
Indeed, the theta transformation is
$\Theta_j(t)=\varepsilon_jt^{-1}\Theta_{7-j}(1/t)$; after changing
$t$ to $1/t$ in the lower Mellin integral, its value at $s=1$ is precisely
the second term in \eqref{seven:L1-formula}.
Because $0<E_1(x)\le e^{-x}/x$, deleting all terms $n>X$ in both sums
has joint absolute error at most
\[
 {\cal T}(c,X)=\frac{2e^{-c(X+1)}}{1-e^{-c}}.
 \tag{S14}\label{seven:tail}
\]
This estimate uses only $|a_n|\le n$ and a geometric series.

The degree-seven field $Q_1$ has
$\zeta_{Q_1}(s)=\prod_{j=0}^6L(s,\psi^j)$.
Here $\psi^0$ means the trivial character of conductor one, whose factor is
$\zeta(s)$, so its Euler factor at seven is retained.
For $j\ne0$ its factors at one are given by the finite identity
\[
 L(1,\psi^j)=
 -\frac{\tau(\psi^j)}{49}
 \sum_{\substack{1\le a<49\\7\nmid a}}
       \overline{\psi^j(a)}
       \log\left(2\sin\frac{\pi a}{49}\right).
 \tag{S15}\label{seven:dirichlet-residue}
\]
This follows by Fourier expansion of $-\log(2\sin\pi x)$ and the
Gauss-sum identity; evenness combines the two Fourier directions.
Pairing conjugate characters makes the resulting residue real and positive.
Finally, the permutation representation for $F_1/\mathbb Q$ is
$(1+\rho)\otimes\bigoplus_{j=0}^6\psi^j$, so Artin formalism gives
\[
 \operatorname*{res}_{s=1}\zeta_{F_1}(s)
 =\operatorname*{res}_{s=1}\zeta_{Q_1}(s)
       \prod_{j=0}^6L_j(1).
 \tag{S16}\label{seven:residue-factorization}
\]
This identity includes every Euler factor, including those at seven.

\subsection{Explicit units and rigorous interval isolation}
\label{seven:certification}

Let
\[
 q(Y)=Y^7-21Y^5-21Y^4+91Y^3+112Y^2-84Y-97,
 \qquad
 \beta=\sum_{a\in\{1,18,19,30,31,48\}}\zeta_{49}^{a}.
 \tag{S17}\label{seven:period}
\]
The set of exponents is the subgroup of order six in
$(\mathbb Z/49\mathbb Z)^\times$. Thus $\mathbb Q(\beta)=Q_1$.
The polynomial $q$ is irreducible and $q(\beta)=0$ by direct arithmetic
modulo $\Phi_{49}$.
Write $\theta=\alpha-\beta$. The polynomial $p$ of degree twenty-one and
the rational function $P/Q$ printed in the finite data below satisfy
\[
 p(\theta)=0,\qquad \alpha=P(\theta)/Q(\theta),\qquad
 \alpha^3=269,\qquad q(\alpha-\theta)=0.
 \tag{S18}\label{seven:field-identity}
\]
Irreducibility of $p$ and $q$, or the relatively prime degrees three and
seven together with these identities, identifies this field as $F_1$.
The automorphism $\tau$ fixes $\alpha$ and is
$\zeta_{49}\mapsto\zeta_{49}^3$ on $Q_1$; its action on $\beta$ is the
degree-six rational polynomial printed below.

The first six units $u_1,\ldots,u_6$ lie in $\mathbb Q(\beta)$ and are
given by their seven coefficients and denominators.
The seventh is
\[
 u_7=-\varepsilon,\qquad
 \varepsilon=270-739\alpha+108\alpha^2.
 \tag{S19}\label{seven:base-unit}
\]
For the explicitly printed $\eta$, the remaining units are
$u_{8+i}=\tau^i\eta$ for $0\le i\le5$.
Exact rational polynomial arithmetic verifies that all thirteen units are
integral and have absolute norm one. One can check integrality either in
the certified maximal order or from the monic integral characteristic
polynomial of multiplication by the element. The discriminant and
signature are
\[
 D_{F_1}=-N^7\,7^{36},\qquad (r_1,r_2)=(7,7).
 \tag{S20}\label{seven:disc-signature}
\]
Independently, $D_{Q_1}=49^6=7^{12}$ by the conductor--discriminant
formula. Its discriminant is coprime to $D_F$, so the product of the
integral bases is a maximal order in the compositum and gives
\eqref{seven:disc-signature}.

Here are complete numerical specifications for the interval calculation.
Generate the integer coefficients in \eqref{seven:coefficients} by a
prime-factor sieve. In \eqref{seven:L1-formula} take $X=14000$ for $j=0$
and $X=550000$ for $j\ne0$, and evaluate all terms by real and complex
ball arithmetic at precision $300$ bits.
The bounds \eqref{seven:tail} are, respectively,
$1.496\cdot10^{-80}$ and $1.328\cdot10^{-62}$.
The resulting enclosures, rounded outward, include
\[
\begin{split}
 L_0(1)&\in
 [\,1.6976757232625543059851766859761711810674
       \ \pm 10^{-40}\,],\\
 \prod_{j=0}^6L_j(1)&\in
 [\,0.6494049966247962459428741324547233673973
       \ \pm 10^{-40}\,].
\end{split}
 \tag{S21}\label{seven:analytic-intervals}
\]
For the regulator, isolate the seven disjoint real roots of $q$ with
complex ball arithmetic, take
$\alpha_0=\exp(\log269/3)$ and
$\alpha_1=\alpha_0(-1+i\sqrt3)/2$, and evaluate the thirteen units at
$(\alpha_0,\beta_j)$ and $(\alpha_1,\beta_j)$.
Form the logarithmic matrix with entries $\log|\sigma(u_i)|$ for real
embeddings and $2\log|\sigma(u_i)|$ for complex pairs, delete one real
row, and take the absolute determinant.
At $2048$ bits this gives the regulator $R_0$ of the displayed unit group
with over $600$ bits of relative accuracy. Arb's enclosures, including
the root isolation and determinant operations, are used throughout
\cite{sevenJohansson2017}.
Rounded outward at lower precision, the relevant values are
\[
\begin{split}
 R_0&\in
 [\,8.7214139234209020456623583408322697739521
       \ \pm10^{-40}\,]\,10^{24},\\
 \operatorname*{res}_{s=1}\zeta_{Q_1}(s)
   &\in[\,0.2996769180948235445839119514335675929714
       \ \pm10^{-40}\,].
\end{split}
 \tag{S22}\label{seven:regulator-intervals}
\]
The positive determinant proves independence of the units.

The analytic class-number formula, in the convention just used, is
\[
 \operatorname*{res}_{s=1}\zeta_k(s)
  =\frac{2^{r_1}(2\pi)^{r_2}h_kR_k}
         {|\mu(k)|\sqrt{|D_k|}}
 \tag{S23}\label{seven:class-formula}
\]
\cite[Theorem~18.13]{sevenSutherland2015}.
If $m_1$ is the index of the displayed thirteen units and $\{\pm1\}$,
\eqref{seven:residue-factorization} and the unrounded ball evaluations give
\[
 \frac{h_{F_1}}{m_1}
 =\frac{\sqrt{N^7\,7^{36}}}
        {2^{13}\pi^7R_0}
   \operatorname*{res}_{s=1}\zeta_{Q_1}(s)\prod_{j=0}^6L_j(1)
 \in[\,7\pm1.11\cdot10^{-59}\,].
 \tag{S24}\label{seven:quotient-interval}
\]
The stated precision and cutoffs, rather than substitution of the rounded
displayed decimals, specify this computation.
Friedman's bound gives
\[
 m_1<B=34885655693683608182649434,
 \qquad
 B\left|h_{F_1}/m_1-7\right|<3.88\cdot10^{-34}<1.
 \tag{S25}\label{seven:index-bound}
\]
The thirteen residue primes and their full evaluation matrix are printed
below. Their determinant is $1$ in $\mathbb F_7$.
Lemma~\ref{seven:isolation} proves $v_7(h_{F_1})=1$.

The same short certificate supplies the base layer without assuming a
class-number search. For \eqref{seven:base-unit}, exact arithmetic gives
$N_{F/\mathbb Q}(\varepsilon)=1$ and
\[
 0<R_\varepsilon=|\log|\varepsilon||<18,\qquad
 \frac{\sqrt N L_0(1)}{2\pi R_\varepsilon}
 \in[\,7\pm4.0\cdot10^{-68}\,].
 \tag{S26}\label{seven:base-quotient}
\]
Thus $m_0=[\mathcal O_F^\times:\langle-1,\varepsilon\rangle]<72$ and
$h_F=7m_0$.
At the prime $491$ the map $\alpha\mapsto444$ gives
$\varepsilon\mapsto88$, while
\[
 138^7\equiv1,\quad138\not\equiv1,\qquad
 88^{70}\equiv223\equiv138^4\not\equiv1\pmod{491}.
 \tag{S27}\label{seven:base-saturation}
\]
This nonzero residue character proves $7\nmid m_0$.
Consequently
\[
 A(F)=A(F_1)=\mathbb Z/7\mathbb Z.
 \tag{S28}\label{seven:finite-inputs}
\]
All integers, unit expressions, coefficient rules, cutoffs and tail bounds
needed for these two assertions have now been specified.

\subsection{Norm control and stabilization}
\label{seven:stability-proof}

The cube classes in $\mathbb F_7$ are $0,1,6$, whereas $269\equiv3$.
Thus $X^3-269$ is irreducible modulo seven, and seven is unramified and inert
in $F$. The local extension at this unique prime is the unramified cubic
extension of $\mathbb Q_7$.
Each cyclotomic extension $Q_n/Q_m$ is totally ramified at seven and
unramified everywhere else. Base change by the unramified cubic local
field preserves this total ramification and its degree. Therefore every
$F_n/F_m$ has exactly one ramified finite prime, which is totally ramified.

\begin{lemma}\label{seven:norm-control}
Let $L/k$ be cyclic of degree $7^r$, with exactly one ramified finite prime,
which is totally ramified. Then the ideal norm induces an isomorphism
\[
 A(L)_{\operatorname{Gal}(L/k)}\ \xrightarrow{\ \sim\ }\ A(k).
 \tag{S29}\label{seven:coinvariant-control}
\]
\end{lemma}

\begin{proof}
Let $H$ be the Hilbert seven-class field of $L$. It is Galois over $k$,
since its defining maximality property is invariant under every
automorphism of $L/k$. Set
$G=\operatorname{Gal}(H/k)$ and
$A=\operatorname{Gal}(H/L)=A(L)$.
An inertia group $I$ above the sole ramified prime has order $7^r$:
$H/L$ is unramified and $L/k$ is totally ramified.
It therefore maps isomorphically to $\operatorname{Gal}(L/k)$, giving
$G=A\rtimes I$.
If $\sigma$ generates $I$, then $[G,G]=(\sigma-1)A$.
Killing $I$ and $[G,G]$ kills all inertia groups and gives the maximal
abelian unramified quotient of $G$, namely $A/(\sigma-1)A$.
The Hilbert seven-class field $H_k$ of $k$ is contained in $H$, since
$H_kL/L$ is unramified, abelian and of seven-power degree. Conversely,
every abelian unramified subextension of $H/k$ lies in $H_k$.
The quotient is therefore $A(k)$.
Under reciprocity this restriction map is the ideal norm, proving the
claim. There is no archimedean inertia of odd order to consider.
\end{proof}

This is the one-prime case of the usual class-field-theoretic norm
calculation; compare
\cite[Corollaries~1.3.9 and 1.3.12]{sevenSharifi}.
The proof includes the finite control map, so that no unidentified
finite kernel is carried into the next argument.

Fix $n\ge1$, put $A_n=A(F_n)$, and regard it as a module over
$R_n=\mathbb Z_7[C_{7^n}]$.
This ring is local with maximal ideal $(7,T)$, because
$R_n/7R_n=\mathbb F_7[T]/(T^{7^n})$.
Lemma~\ref{seven:norm-control} gives
\[
 A_n/TA_n\simeq A(F),\qquad
 A_n/\omega_1(T)A_n\simeq A(F_1),\qquad
 \omega_1(T)=(1+T)^7-1=T\nu(T),
 \tag{S30}\label{seven:two-quotients}
\]
where
$\nu(T)=7+21T+35T^2+35T^3+21T^4+7T^5+T^6$.
For $n=1$ the second quotient is simply $A_1$, so the statement remains
valid. Both quotients have order seven by \eqref{seven:finite-inputs}.
Since $\omega_1A_n\subseteq TA_n$, equality of their indices implies
\[
 TA_n=\omega_1A_n=\nu(T)\,TA_n.
\]
But $\nu(T)\in(7,T)$, so Nakayama's lemma applied to the finite
$R_n$-module $TA_n$ gives $TA_n=0$.
The first quotient in \eqref{seven:two-quotients} now proves
\eqref{seven:class-stability}.
Applying Lemma~\ref{seven:norm-control} to every $F_n/F_m$ shows that all
norms are surjective, hence isomorphisms between groups of order seven.
The resulting inverse limit is $\Lambda/(7,T)$.
This also proves directly the needed case of Fukuda's theorem
\cite[Theorem~1(1)]{sevenFukuda1994}.

The same norm lemma applied to $Q_n/\mathbb Q$, followed by Nakayama,
gives
\[
 A(Q_n)=0 \qquad(n\ge0).
 \tag{S31}\label{seven:rational-layer}
\]
Indeed its coinvariants are $A(\mathbb Q)=0$.

\subsection{The two faithful cubic-character components}
\label{seven:character-components}

Let $c$ be complex conjugation. Then
$\operatorname{Gal}(E_n/Q_n)=S_3$ with
$c\delta c=\delta^{-1}$, and $E_n^{\langle c\rangle}=F_n$.
Set $C_n=A(E_n)$.
For an extension of degree $d$ prime to seven, extension and norm of ideals
satisfy
$Nj=d$ and $jN=\sum_g g$.
Since $d$ acts invertibly on seven-primary groups, extension identifies the
lower class group with the invariant subgroup, without a ramification
restriction. Thus
\[
 C_n^{c=1}\simeq A(F_n),\qquad
 C_n^\Delta\simeq A(K_n),\qquad
 (C_n^\Delta)^{c=1}\simeq A(Q_n)=0.
 \tag{S32}\label{seven:prime-to-seven-descent}
\]
For $\varphi\in\{1,\chi,\chi^{-1}\}$ put
$e_\varphi=\frac13\sum_{d\in\Delta}\varphi(d)^{-1}d$.
The idempotents split
$C_n=C_n^1\oplus C_n^\chi\oplus C_n^{\chi^{-1}}$, and $c$ exchanges the
last two summands.
Since the $c$-fixed part of $C_n^1$ vanishes, projection gives an
isomorphism
\[
 e_\chi:C_n^{c=1}\ \xrightarrow{\ \sim\ }\ C_n^\chi,
 \qquad x\longmapsto e_\chi x,
 \qquad y\longmapsto y+cy
 \tag{S33}\label{seven:projection}
\]
with the displayed inverse.
Extension of ideals, the idempotents, and these projections commute with
the cyclotomic action and with the norm maps in the compositum tower.
Equations \eqref{seven:class-stability} and \eqref{seven:projection}
therefore prove \eqref{seven:iwasawa} for both faithful characters.

Only $(C_n^1)^{c=1}=0$ was used; deleting $C_n^1$ itself would be false.
For example, $K$ has class number one by its Minkowski bound
$(2/\pi)\sqrt3<2$. In $K_n/K$ the two primes above seven are totally
ramified, and each unit of $K$ is the norm of a unit, since the norm on
$\mu_6$ is raising to $7^n$.
The cyclic ambiguous class-number formula consequently gives
$|A(K_n)^{\operatorname{Gal}(K_n/K)}|=7^n$.
This growth lies in the conjugation-odd part and does not affect
\eqref{seven:projection}; the relevant cyclic class-field exact sequence
is given in \cite[Theorem~1.3.16]{sevenSharifi}.

\subsection{Strict local control and specialization}
\label{seven:control}

Seven splits into two primes in $K$, and $E/K$ is unramified and inert of
degree three at each of them. Hence each prime over seven in $E_n$ is fixed
by $\Delta$.
If
\[
 C_n'=\operatorname{Cl}(\mathcal O_{E_n}[1/7])\{7\},
 \qquad M'=\varprojlim_n C_n',
\]
the kernel of $C_n\to C_n'$ is generated by the seven-prime classes and
is $\Delta$-trivial. It follows that
\[
 (M')^\chi
 \simeq\left(\varprojlim_n C_n\right)^\chi
 \simeq\Lambda/(7,T).
 \tag{S34}\label{seven:split-module}
\]
Class field theory identifies $M'$ with the maximal unramified abelian
pro-seven Galois group over $E_\infty$ in which every seven-prime splits
completely. This is the norm inverse limit of the finite Hilbert
$S$-class-field correspondences
\cite[Remark~1.3.21 and Proposition~3.2.20]{sevenSharifi}.

Here one can check the global and local descent maps explicitly.
Let $G=\Delta\times\Gamma$ and $D=D_{\chi,\kappa}$.
The element $\chi(\delta)-1$ is a unit in $\mathcal O$, so $D^\Delta=0$.
Cohomology in positive degree for the order-three group $\Delta$ vanishes
on seven-primary modules, by averaging.
The Hochschild--Serre sequence therefore gives
$H^i(G,D)=0$ for all $i\ge0$.
The inflation--restriction sequence is consequently an isomorphism
\[
 H^1(K,D)\ \xrightarrow{\ \sim\ }\
 H^1(E_\infty,D)^G
 =\operatorname{Hom}_{\rm cts}(G_{E_\infty},D)^G.
 \tag{S35}\label{seven:global-inflation}
\]
At either seven-prime the decomposition group in $G$ is all of
$\Delta\times\Gamma$: the local extension is the compositum of the
unramified cubic extension and the totally ramified cyclotomic tower.
The identical argument shows that local inflation--restriction is an
isomorphism there. Thus a class vanishes on the full local decomposition
group before descent if and only if its descended homomorphism does.
At a finite prime away from seven the inertia quotient in $G$ has order
one or three; its positive cohomology again vanishes. This makes
restriction injective on the inertia cohomology, so the unramified
condition is also equivalent before and after descent.
These are precisely the two local conditions in
\eqref{seven:strict-definition}. We obtain
\[
 \operatorname{Sel}_{\rm str}(K,D)
 =\operatorname{Hom}_{\rm cts}(M',D)^{\Delta\times\Gamma}.
 \tag{S36}\label{seven:Hom-description}
\]

For clarity about character and inverse conventions, equivariance means
\[
 f(\delta x)=\chi(\delta)f(x),\qquad f(\gamma x)=u f(x).
 \tag{S37}\label{seven:orientation}
\]
The first condition selects the natural $\chi$-component of $M'$.
On that component $\gamma=1$ and $7=0$ by \eqref{seven:split-module}.
Hence \eqref{seven:Hom-description} is $D[7,u-1]$, which, since
$\mathcal O$ is a discrete valuation ring, is noncanonically
$\mathcal O/(7,u-1)$.
This proves \eqref{seven:all-characters}.
Writing a contragredient convention instead replaces $u-1$ by
$u^{-1}-1=-u^{-1}(u-1)$, the same ideal. No Tate twist or unspecified
specialization kernel occurs in this argument.

\subsection{Weights, wild characters, and coefficient normalization}
\label{seven:weights}

For $e\in\mathcal O$, the convergent exponential gives the continuous
character determined by
$\kappa(\gamma)=\exp(-e\log8)$.
It satisfies $u-1\in7\mathcal O$. In particular this includes
$\kappa=\langle\varepsilon_7\rangle^{-e}$ for $e\in\mathbb Z_7$, including
$e=0$, and
\[
 \operatorname{Sel}_{\rm str}
   \bigl(K,(V/\mathcal O)(\chi\langle\varepsilon_7\rangle^{-e})\bigr)
 \simeq\mathcal O/7\mathcal O.
 \tag{S38}\label{seven:principal-weights}
\]
Over $\mathbb Z_7$ every such group has length one.
Over a ramified coefficient ring its $\mathcal O$-length is
$e(V/\mathbb Q_7)$; it must not be replaced by a length-one quotient by
a uniformizer.

More generally, suppose $\mathcal O$ contains a primitive $7^r$th root
$\zeta$, where $r\ge1$, and put
$u=\zeta\exp(-e\log8)$ with $e\in\mathcal O$.
Since
$v_7(\zeta-1)=1/(6\cdot7^{r-1})<1$
and $\exp(-e\log8)-1\in7\mathcal O$, the two summands in
$u-1=(\zeta-1)+\zeta(\exp(-e\log8)-1)$ have different valuations.
Thus
\[
 \operatorname{Sel}_{\rm str}(K,D_{\chi,\kappa})
 \simeq\mathcal O/(\zeta-1),\qquad
 \frac{\log\#\operatorname{Sel}_{\rm str}(K,D_{\chi,\kappa})}
      {[V:\mathbb Q_7]}
 =\frac{\log7}{6\cdot7^{r-1}}.
 \tag{S39}\label{seven:wild-cost}
\]
The valuation formula follows by applying the Eisenstein polynomial
$\Phi_{7^r}(1+X)$ to $\zeta-1$.
For the principal branch \eqref{seven:principal-weights}, the same
normalized logarithmic size is $\log7$.
For an arbitrary continuous $\kappa$, the completely general answer
remains \eqref{seven:all-characters}; no parametrization of all such
characters by $\exp(-e\log8)$ is assumed.

Finally consider a coefficient family written in the original finite
character convention as
\[
 D_{j,a,e}=(V/\mathcal O)
    \bigl(\chi^j\xi_7^a\varepsilon_7^{-e}\bigr),\qquad
 j\in\{1,2\},\quad0\le a<6,\quad e\ge1,
 \qquad \xi_7=\omega_7^v,\quad v\in\{1,5\}.
 \tag{S40}\label{seven:matched-coefficients}
\]
If the finite residual characters are matched, then
$\xi_7^a\omega_7^{-e}=1$ exactly, because reduction is injective on
$\mu_6$. Equivalently, $va\equiv e\pmod6$, which determines one $a$
for every $e$. Thus each matched coefficient is exactly
$(V/\mathcal O)(\chi^j\langle\varepsilon_7\rangle^{-e})$, and
\[
 s(e)=1\quad\hbox{for every positive matched weight},\qquad
 \sum_{j=1}^{2}\sum_{e=1}^{N}
 \frac{\log\#\operatorname{Sel}_{\rm str}(K,D_{j,a(e),e})}
      {[V:\mathbb Q_7]}
 =2N\log7.
 \tag{S41}\label{seven:matched-cost}
\]
There is no exceptional weight in this matched family. Statements about
other tame characters or additional factors require their own local
matching and are not used here.

\subsection{Complete finite unit and residue data}
\label{seven:finite-data}

The following arrays are exact integers. Coefficients for $p$, $q$, $P$ and
$Q$ are in descending order. The coefficients for $\tau$, the six small
units, and the cubic unit are in ascending order. A small-unit numerator
$[a_0,\ldots,a_6]$ represents $\sum a_i\beta^i$; the cubic-unit numerator
represents a polynomial in $\alpha$.
The twenty-one coefficients of $\eta$ are ordered with $\alpha$ varying
first: coefficient number $3j+i+1$ multiplies $\alpha^i\beta^j$,
$0\le i\le2$, $0\le j\le6$.
To keep every integer on the page, each coefficient of $\eta$ is printed as
$[s,a_0,a_1,a_2]$ in base $B_*=10^{100}$, meaning
$s(a_0+B_*a_1+B_*^2a_2)$; no leading zeroes are required.

Each residue triple is $(\ell,r,z)$: reduce $\theta$ to $r$ modulo $\ell$
and use $z$ in \eqref{seven:residue-character}. The thirteen primes with
these residues are distinct, even when their rational primes coincide.
The roots of $p$ are simple at all these residues; the derivative residues,
in the same order, are
\[
 153,14,99,174,74,74,436,46,119,368,219,76,606.
\]
Hensel's lemma therefore supplies the corresponding maps of the full
integer ring to $\mathbb F_\ell$.
The rows of the final matrix are in this order and the columns are
$u_1,\ldots,u_{13}$.
Polynomial substitution, the norm and integrality tests, the finite-field
powers in \eqref{seven:residue-character}, and Gaussian elimination modulo
seven are a complete exact verification of these finite witnesses.

\begingroup
\fontsize{7.5}{9.5}\selectfont
\begin{verbatim}
p = [1, 0, -63, -1820, 1596, 47859, 1550766, -771102, 74382, -773205433, 265573350, -12672069060,
  225337203987, -164087914221, 5350549479537, -37293576067279, 55623366413664, -844644350170257,
  3089185218698537, -5358907459195815, 42933560367107697, -74554880595863287]
q = [1, 0, -21, -21, 91, 112, -84, -97]
P = [7, 0, -252, 43540, 3115, -767592, 28694829, 4854766, -154925106, 3054427187, 508674397, 875124348,
  32367791232, 6622577365]
Q = [28, 0, -819, 72163, 8253, -926394, 25924017, 4555698, -71871954, 1423159430, 242638326,
  1061561823, 4574293518]
tau_numerator = [3210, -2025, -2700, 820, 495, -51, -15]
tau_denominator = 589
small_unit_numerators = [
  [48, 18, -65, -13, 17, 2, -1],
  [-842, 2689, 5156, -305, -1495, -78, 81],
  [-5185, -1631, 4697, 659, -1028, -69, 49],
  [5927, -5217, -7545, 1048, 1825, 44, -91],
  [5085, -3117, -2389, 743, 330, -34, -10],
  [10258, -972, -10131, -431, 2358, 164, -125],
]
small_unit_denominators = [1, 589, 589, 589, 589, 589]
cubic_unit_numerator = [-270, 739, -108]
eta_denominator = 1767
eta_coefficients_in_base_10_power_100 = [
  [1,
   8110910039399052799388839509913995719324441857259887672553997475894562301601700471535299521506405517,
   3431137402558321943275234065603513747732585590776418713813607417676553074824870235139769709481053259,
   3940696968887],
  [1,
   2863301462649584998067112611608264044687998196095409917598114508277713738730199925757154728888953055,
   970812101337790668521585581275454765494152664299295398136805535577832483461611290913002007827311071,
   610457752147],
  [1,
   4821724478914456225350364309117571242884298504341022610343187741267858848133670184073817508713496298,
   3511990174213849653115720097913645282815157493086558732902294533502899394979501897922579875243606017,
   94566689623],
  [1,
   760344159491231074673603899303458724681328958152289201635437873131003916558596294478134311589863496,
   2317480636152653863897711154536450131210677689030596417135280558908466878562464587662862833652562118,
   7059069004615],
  [1,
   5713646533666935748952660526437906230094439209439484821582261626527132028034044010356200033729959598,
   8155918751592805126188738962782659016016516467436799064429460997991815680374480211693583033342107571,
   1093528233921],
  [1,
   5458376416976996449886152263319535978441113085075418487179464286576609680040631434748810087394632375,
   2815314908475193445673806280168229698608681052505411231685686073970229478053287079417408542790496382,
   169399675453],
  [1,
   9104803622062202912042424298875099840356676073442457598635399038039260264973760983790173980634551253,
   5209978760288616211810004718597728762354940779407360157403598658202633153414464523403435572630242052,
   1981995488635],
  [1,
   3229190277068481099567369991097155293895283837272933560810654867460601672084248677453637066047248351,
   1066602590408197039412448163000624315584233985107205159080701191411246693038208294958734638546267908,
   307033126452],
  [1,
   4719833413501908223670674319560071208620416677277259622958912093004163789224972446770714978198517369,
   3589225343881354158630968166458173644607882286774087545514248915094188084702046502205303007048245172,
   47562843245],
  [-1,
   7252159511744921579747663915866179139747715077640636824080064020836768547568447722702965774255379864,
   2117795437844894019014366078716012207469677439109909489100232555667366018470798129278480577073347195,
   1862911672364],
  [-1,
   4596952397808038061504443355452507844075138976341878918188397033987218237849921882684029382396888813,
   5548014240268422805702685958606065400771312345386887361270938503247549652804788711898670621959818004,
   288585719972],
  [-1,
   8576573868236387713297262546743178973121352627702349510026823016588229660225220810296942340315536894,
   7246874522795187239999206975901393322813195843039386133915845845849949086654946548202311744607548915,
   44705134981],
  [-1,
   7281143800513613602674102889822350672969769041691650309307424554802016538032032999848411121384589417,
   7247656128371841488568688653380385878443917715304128569138890771204319093551264064535444911867233869,
   870696576112],
  [-1,
   2603158096105930003776498071993161225442241308288062551734654529270377694879518302937682334035359418,
   8801224015145560570571877684554110083953909461438848969946722322537320157416843524148670087474475623,
   134880575403],
  [-1,
   2144340361888752919586412430689254661717183112022793764572101308645944000985194209283484734442009651,
   4344877732372528348076072331962740358969404598366277209320255591698654981676044669701093523564674725,
   20894500013],
  [1,
   1658789263126353124561691351357597966888983262121120607405319802389834239184987248275993385042870297,
   4563350942981641714892768486786517151934762869785113709687288930688939203417092386426646047883401037,
   47445317587],
  [1,
   3334821144767026713097794436460203379080665880747592876486755438276015329827178694499770466419371729,
   4734315972881480902169887703698274649856456832930081114001007355628760067288435055252831111892960802,
   7349806938],
  [1,
   2179355790946349106219951244546405670330306327971639576268515952647904221008259146456375776668744595,
   7511295224937829629681218123710093110445808535628505578790220369499924491111462384032633379660007227,
   1138566770],
  [1,
   1940103518950861831403240058854220493243755110925587464386601540881397210941906392843448161215669392,
   7398981907903882144917187999221062721013600253004142759479698501891917306885982478145740682084331734,
   43903316832],
  [1,
   3471460889921431092203467331457652901571380143073529228472752287967796312304491085909952970479760717,
   4781793234571883106444295805054726784965688101009565012732394759877052590133372264885207805147053091,
   6801111660],
  [1,
   1294025870817364042382316796528664510562816139774755081398276045466303764150138262507798612436708108,
   7200960424005516814554303735541605542146140398242153380225569364053345312160643463981887310532844365,
   1053567774],
]
prime_root_zeta_triples = [
  [197, 164, 104],
  [197, 119, 104],
  [197, 110, 104],
  [197, 69, 104],
  [197, 50, 104],
  [197, 29, 104],
  [491, 435, 138],
  [491, 431, 138],
  [491, 411, 138],
  [491, 107, 138],
  [491, 98, 138],
  [491, 93, 138],
  [883, 874, 707],
]
residue_character_matrix = [
  [2, 3, 4, 2, 6, 5, 0, 5, 5, 2, 3, 3, 5],
  [5, 6, 2, 5, 1, 1, 0, 5, 5, 5, 2, 3, 3],
  [0, 1, 3, 1, 0, 2, 0, 5, 5, 5, 5, 2, 3],
  [6, 0, 6, 2, 5, 5, 0, 3, 5, 5, 5, 5, 2],
  [5, 4, 0, 0, 2, 6, 0, 2, 3, 3, 5, 5, 5],
  [2, 5, 1, 5, 4, 0, 0, 3, 3, 5, 5, 5, 5],
  [6, 6, 2, 0, 4, 5, 4, 3, 4, 1, 0, 6, 4],
  [4, 4, 2, 5, 0, 3, 4, 3, 3, 4, 1, 0, 6],
  [5, 2, 0, 3, 2, 0, 4, 1, 0, 6, 4, 3, 3],
  [3, 0, 6, 3, 0, 6, 4, 4, 3, 3, 4, 1, 0],
  [3, 2, 0, 0, 6, 0, 4, 4, 1, 0, 6, 4, 3],
  [0, 0, 4, 6, 0, 4, 4, 6, 4, 3, 3, 4, 1],
  [6, 4, 2, 5, 6, 6, 1, 1, 4, 3, 4, 6, 3],
]
determinant_modulo_7 = 1
\end{verbatim}
\endgroup
